\chapter{Adjoint Sensitivity in Practice}

\begin{flushright}
\textit{Reverse the calculation, and one output can speak about many inputs.}
\end{flushright}

\bigskip

You need the sensitivities of a complex disease model with 30~parameters.
The model is your own code, you wrote it yourself, so it is differentiable.
You know from Chapter~6 that the adjoint requires only 2~ODE solves instead
of 30~FSE solves.

But deriving the adjoint equations by hand for 30~coupled nonlinear ODEs would
take weeks and would almost certainly introduce errors in the derivation.

Is there a way to obtain the adjoint computation without deriving the adjoint
equations by hand? Yes, algorithmic differentiation does this automatically~\cite{naumann2012art}.
And there is a way to use it from MATLAB today, without writing
a single AD routine from scratch.

%%------------------------------------------------------------
\section{The Annuity Function: A Worked Comparison}
\label{sec:annuity}
%%------------------------------------------------------------

Before examining the computational tools, we work through a concrete
example where all three methods, forward mode, reverse mode,
and symbolic differentiation, can be compared exactly.

\subsection{The Problem}

An annuity pays a fixed amount $P$ per period for $n$~periods at
interest rate $r$. The present value is
\begin{equation}
  A(P, r, n) = P\,\frac{1 - (1+r)^{-n}}{r}.
  \label{eq:annuity}
\end{equation}
We want all three sensitivity indices: $\SI{P}{A}$, $\SI{r}{A}$,
$\SI{n}{A}$.

\subsection{Forward Mode: Three Passes}

In forward mode we differentiate the computation one parameter at a time,
propagating $\dot{v}_k = \partial v_k/\partial q$ through intermediate
values. We trace the computation of $A$ through the following steps:
\[
  v_1 = (1+r)^{-n}, \quad
  v_2 = 1 - v_1, \quad
  v_3 = v_2/r, \quad
  A = P v_3.
\]

For $q = P$: set $\dot{P} = 1$, $\dot{r} = 0$, $\dot{n} = 0$.
\[
  \dot{v}_1 = 0, \quad \dot{v}_2 = 0, \quad \dot{v}_3 = 0, \quad
  \dot{A} = \dot{P}\,v_3 + P\,\dot{v}_3 = v_3.
\]
So $\partial A/\partial P = v_3 = (1-(1+r)^{-n})/r$.

For $q = r$: set $\dot{r} = 1$, $\dot{P} = 0$, $\dot{n} = 0$.
\[
  \dot{v}_1 = n(1+r)^{-n-1}, \quad
  \dot{v}_2 = -n(1+r)^{-n-1}, \quad
  \dot{v}_3 = \frac{\dot{v}_2\,r - v_2}{r^2}, \quad
  \dot{A} = P\,\dot{v}_3.
\]

For $q = n$: set $\dot{n} = 1$, $\dot{P} = 0$, $\dot{r} = 0$.
\[
  \dot{v}_1 = -\ln(1+r)\,(1+r)^{-n}, \quad
  \dot{v}_2 = \ln(1+r)\,(1+r)^{-n}, \quad
  \dot{v}_3 = \dot{v}_2/r, \quad
  \dot{A} = P\,\dot{v}_3.
\]

Three passes through the computation graph give three derivatives.
Each pass involves the same number of operations as one evaluation
of $A$, so the total cost is approximately $3\times$ the cost
of a single function evaluation.

\subsection{Reverse Mode: One Pass}

In reverse mode, we compute all three derivatives in one backward
pass. Set $\bar{A} = 1$. Work backward through the computation:
\begin{align*}
  \bar{P}  &= \bar{A}\,v_3 = v_3, \\
  \bar{v}_3 &= \bar{A}\,P = P, \\
  \bar{v}_2 &= \bar{v}_3 / r, \\
  \bar{r}   &= \bar{v}_3 \cdot (-v_2/r^2) + \bar{v}_1\cdot\pd{v_1}{r}, \\
  \bar{v}_1 &= -\bar{v}_2, \\
  \bar{n}   &= \bar{v}_1 \cdot \pd{v_1}{n}
             = -\bar{v}_2 \cdot \bigl(-\ln(1+r)(1+r)^{-n}\bigr).
\end{align*}

One backward pass gives $\partial A/\partial P$, $\partial A/\partial r$,
and $\partial A/\partial n$ simultaneously. Total cost: approximately
$2\times$ one function evaluation (one forward pass to store values;
one backward pass to accumulate derivatives), independent of
how many parameters $m$ there are.

\subsection{Cost Summary}

\subsection{Symbolic Differentiation}

As a third comparison, we differentiate the formula~\eqref{eq:annuity}
by hand (or by a computer algebra system). Define $v = (1+r)^{-n}$.
Then $A = P(1-v)/r$, giving:
\[
  \pd{A}{P} = \frac{1-v}{r}, \quad
  \pd{A}{r} = P\!\left(\frac{nv/(1+r) - (1-v)/r}{r}\right),
  \qquad
  \pd{A}{n} = P\,\frac{\ln(1+r)\,v}{r}.
\]
Symbolic differentiation produces exact closed-form expressions that can
be evaluated at any $(P, r, n)$ with a single function evaluation;
no passes through the computation graph are required at query time.
The cost is paid upfront in the derivation; the payoff is that evaluation
is cheap thereafter.

\subsection{Exact Operation Counts}

For the annuity function, we can count arithmetic operations precisely.
The function $A$ requires: 1 addition, 1 exponentiation, 1 subtraction,
2 divisions, 1 multiplication = 6 operations.

Forward mode for one parameter: 6 operations for the primal values,
plus 6 tangent operations = 12 total. For $m = 3$ parameters: 36 operations.

Reverse mode: 12 operations (6 forward + 6 backward) regardless of $m$.
For $m = 3$: still 12 operations.

The ratio: forward mode costs $\approx 2m$ times a single evaluation;
reverse mode costs $\approx 2$ times a single evaluation. For the annuity
with $m = 3$, forward mode uses $3\times$ more operations than reverse mode.
For a climate model with $m = 500$, forward mode uses $250\times$ more.
For a neural network with $m = 10^7$ parameters and one scalar loss function,
reverse mode still requires only one backward pass --- the same cost as the
annuity --- while forward mode would require $10^7$ separate sweeps.
This is why reverse-mode AD made modern deep learning computationally feasible.

\begin{center}
\begin{tabular}{llll}
\toprule
Method & Passes & Operations (annuity, $m=3$) & Scales as \\
\midrule
Forward mode & $m = 3$ & $\approx 36$ & $O(m)$ \\
Reverse mode & 2 & $\approx 12$ & $O(1)$ in $m$ \\
Symbolic diff & 1 (at query time) & $\approx 6$ & $O(1)$ in $m$ \\
Finite diff & $2m + 1 = 7$ evals & $\approx 42$ & $O(m)$, approx \\
\bottomrule
\end{tabular}
\end{center}

At nominal values $P = 1000$, $r = 0.05$, $n = 20$, all methods give:
$\SI{P}{A} = 1$, $\SI{r}{A} \approx -0.604$, $\SI{n}{A} \approx 0.482$.
(The derivation of the analytic formulas is Exercise~7.1.)

\begin{selfcheck}
From the annuity formula~\eqref{eq:annuity}, compute $\SI{P}{A}$ analytically
using the derivative form $\SI{P}{A} = (P/A)(\partial A/\partial P)$.
Verify that it equals $1$ for any nominal values of $P$, $r$, $n$.
\end{selfcheck}
\begin{scAnswer}
$\partial A/\partial P = (1-(1+r)^{-n})/r = A/P$.
So $\SI{P}{A} = (P/A)(A/P) = 1$.

This holds for any $r$, $n$: the present value is linear in $P$,
so a $1\%$ increase in the payment amount produces exactly a $1\%$
increase in present value, regardless of the interest rate or
number of periods.
\end{scAnswer}

%%------------------------------------------------------------
\section{Algorithmic Differentiation}
\label{sec:AD}
%%------------------------------------------------------------

Algorithmic differentiation (AD), also called automatic
differentiation, is the chain rule applied systematically
to every arithmetic operation in a computer program.
It is neither symbolic differentiation (which manipulates formulas)
nor numerical differentiation (which approximates derivatives).
It is exact to machine precision and applies to any differentiable code.

Students sometimes expect AD to produce a formula for the derivative,
like a computer algebra system would. It does not.
AD produces the numerical value of the derivative at a specific
input, by accumulating the chain rule through the sequence of
arithmetic operations the program performs.

Similarly, reverse-mode AD is not the same as numerical differentiation.
A centered-difference estimate of $\partial A/\partial r$ has error
$O(h^2)$ from the truncation of the Taylor series; reverse-mode AD
gives the exact derivative (to floating-point precision) with no
step-size selection required.

\subsection{Forward Mode}

Forward-mode AD attaches a tangent value $\dot{v}_k = \partial v_k/\partial q$
to each intermediate variable $v_k$, for one chosen parameter $q$.
At each arithmetic operation, the tangent is updated using the chain rule:
\[
  v_k = f(v_i, v_j) \implies
  \dot{v}_k = \pd{f}{v_i}\dot{v}_i + \pd{f}{v_j}\dot{v}_j.
\]
At the end of the computation, $\dot{v}_{\text{output}}$ contains
$\partial J/\partial q$. Repeat for each parameter: $m$~passes total.

\subsection{Reverse Mode}

Reverse-mode AD attaches an adjoint value $\bar{v}_k = \partial J/\partial v_k$
to each intermediate variable. In the backward pass, these are accumulated
using the chain rule in reverse:
\[
  v_k = f(v_i, v_j) \implies
  \bar{v}_i \mathrel{+}= \pd{f}{v_i}\bar{v}_k, \quad
  \bar{v}_j \mathrel{+}= \pd{f}{v_j}\bar{v}_k.
\]
The forward pass computes and stores all intermediate values.
The backward pass accumulates all $\partial J/\partial q_i$ simultaneously.
Total cost: one forward pass + one backward pass, regardless of $m$.

Reverse-mode AD is the discrete-computation-graph version
of the adjoint ODE derived in Chapter~6. The adjoint ODE propagates
$\vec{\lambda}(t)$ backward through a continuous system; reverse-mode
AD propagates $\bar{v}_k$ backward through a discrete computation graph.
The underlying mathematical operation is identical: transposition of
the derivative operator.

\unifyingidea{2} Reverse-mode AD implements Unifying Idea~2 automatically:
the adjoint is the transpose in whatever space you are working in, and AD
constructs that transpose through the backward pass without requiring the analyst to derive the adjoint equations by hand.
For a model written as differentiable code, this eliminates the main
practical barrier to adjoint SA: instead of deriving and implementing
the adjoint ODE, the analyst uses an AD library that differentiates
the code directly. The adjoint equation is implicit in the
reverse accumulation; the analyst never writes it explicitly.

\subsection{Available Tools}\index{algorithmic differentiation!software}

The 2009 source material for this book referenced ADIFOR, DAFOR,
and TAMC, tools that are now outdated or unmaintained.
Table~\ref{tab:ADtools} lists the current ecosystem.
For a comprehensive survey of AD methods, see Baydin et al.~\cite{baydin2018automatic}
and Margossian~\cite{margossian2019review}.

\begin{table}[htb]
\centering
\renewcommand{\arraystretch}{1.2}
\begin{tabular}{p{2cm}p{4.2cm}p{2.5cm}p{3.5cm}}
\toprule
Language & Tool & Mode(s) & Notes \\
\midrule
MATLAB & Symbolic Math Toolbox & Symbolic & Exact; limited scalability \\
MATLAB & Deep Learning Toolbox \texttt{dlarray} & Reverse & AD for code; GPU support \\
Python & JAX & Forward + Reverse & NumPy-compatible; JIT compilation \\
Python & PyTorch \texttt{autograd} & Reverse & Widely used in ML \\
Julia & ForwardDiff.jl & Forward & High performance \\
Julia & Zygote.jl & Reverse & Source-to-source; flexible \\
C/Fortran & ADOL-C & Forward + Reverse & Mature; still active \\
C/Fortran & Tapenade & Forward + Reverse & Source transformation \\
\bottomrule
\end{tabular}
\caption{Current AD tools by language (as of 2025).
         For MATLAB users, the Deep Learning Toolbox \texttt{dlarray}
         provides reverse-mode AD for arbitrary differentiable code.
         For large-scale scientific computing, JAX (Python) or ADOL-C
         (C/Fortran) are the most commonly used options.}
\label{tab:ADtools}
\end{table}

\begin{selfcheck}
Explain in one sentence each why: (a)~AD is not symbolic differentiation;
(b)~reverse-mode AD is not the same as centered finite differences;
(c)~forward-mode AD requires $m$ passes while reverse-mode requires 2.
\end{selfcheck}
\begin{scAnswer}
(a) Symbolic differentiation manipulates formulas algebraically;
AD applies the chain rule numerically to the sequence of operations
in the code, producing a number rather than a formula.

(b) Centered finite differences approximate the derivative with error
$O(h^2)$ via $(f(p+h)-f(p-h))/(2h)$; reverse-mode AD computes the
exact derivative (to floating-point precision) by accumulating the
chain rule backward through the computation graph.

(c) Forward mode propagates the tangent $\dot{v}_k = \partial v_k/\partial q$
for one parameter $q$ per pass, requiring one pass per parameter.
Reverse mode propagates the adjoint $\bar{v}_k = \partial J/\partial v_k$
for one response $J$ per backward pass, giving all $m$ derivatives simultaneously.
\end{scAnswer}

%%------------------------------------------------------------
\section{The Master Decision Table}
\label{sec:mastertable}
%%------------------------------------------------------------

With all methods now in hand, Table~\ref{tab:masterSA} consolidates
the decision rules developed across Chapters~3--7.

\begin{table}[htb]
\centering
\footnotesize
\renewcommand{\arraystretch}{1.2}
\begin{tabular}{p{2.8cm}p{1.9cm}p{2.4cm}p{1.8cm}p{3.5cm}}
\toprule
Model type & $m$ vs.\ $r$ & Method & Accuracy & Cost \\
\midrule
Black-box & Any & FD (Ch.~3) & $O(h^2)$ & $2m{+}1$ evals \\
Analytic, simple & $m \leq r$ & FSE (Ch.~4) & Exact & $m$ solves \\
Analytic, diff.\ code & $m \gg r$ & AD, reverse & Exact & 2 passes \\
Analytic, manual & $m \gg r$ & Adjoint (Ch.~6) & Exact & 2 solves \\
Correlated POIs & Any & Reparameterize & --- & Depends \\
Near bifurcation & Any & FSE + caution & Exact* & $m$ solves \\
Large uncertainty & Any & Global SA (Ch.~9) & Variance & $N(m{+}2)$ evals \\
\bottomrule
\end{tabular}
\caption{Master decision table for SA method selection;
         see Pianosi et al.~\cite{pianosi2016sensitivity} for a broader SA survey.
         The primary driver is whether the model is black-box or
         differentiable, and whether $m \gg r$ or $m \leq r$.
         (*Near a bifurcation, the FSE is exact but the SI diverges;
         use the result with caution and consult Chapter~8.)
         \textbf{AD vs analytic adjoint:}
         use reverse-mode AD when the model is implemented in differentiable code
         and memory allows storing the computation tape.
         Use the analytic adjoint (Chapter~6) when the model is a nonlinear ODE or PDE,
         memory is constrained, or the numerical integrator (e.g., adaptive \texttt{ode45})
         is not itself differentiable.}
\label{tab:masterSA}
\end{table}

There is no single universally correct method. The right choice depends
on the model, the available tools, and the question being asked.
A useful heuristic: start with finite differences (Chapter~3) to
get SA results quickly, then switch to the adjoint (Chapter~6 or
AD) if the computation time is unacceptable or if $m$ is large.

One principle applies universally: whatever method you use, always
sanity-check at least two sensitivity indices against an independent
method before trusting the full result. For finite differences,
verify one analytic SI if available. For the adjoint, verify two
values against finite differences. This cross-check catches the
majority of implementation errors.

You should not need to choose between finite differences and the adjoint.
Both serve distinct roles: FD is the right tool for black-box models;
the adjoint is the right tool for white-box (differentiable) models
with many parameters. They are complementary, not competing.
But each tool has failure modes that practitioners encounter repeatedly.

%%------------------------------------------------------------
\section{Practical Guidance and Common Pitfalls}
\label{sec:pitfalls}
%%------------------------------------------------------------

\subsection{Discretize-Then-Adjoint\index{discretize-then-adjoint} vs.\ Adjoint-Then-Discretize}

When computing adjoint sensitivities for ODE systems, two strategies
are available. In the \textbf{adjoint-then-discretize} approach
(used throughout Chapter~6), the adjoint ODE is derived analytically
from the continuous ODE and then solved numerically. In the
\textbf{discretize-then-adjoint} approach, the ODE is first
discretized into a finite-difference recurrence, and the adjoint
is then derived for that discrete system.

For linear ODEs with fixed step sizes, both approaches give identical
results. For nonlinear ODEs with adaptive solvers, they can differ:
the adjoint-then-discretize approach may give sensitivities that
disagree with finite-difference verification because the continuous
adjoint assumes a smooth forward trajectory that the adaptive solver
does not produce. The discretize-then-adjoint approach is always
consistent with finite differences but requires differentiating
through the step-size control logic, which is non-trivial.

For the applications in this book, the adjoint-then-discretize approach
with a fixed-step-size forward solver is the most practical choice.
Sirkes and Tziperman~\cite{sirkes1997finite} discuss the implications
for geophysical models where both approaches are used in practice.

\subsection{Memory Requirements for the Adjoint}\index{checkpointing}

The adjoint ODE requires the forward solution $\vec{u}(t)$ stored at
every time point where the adjoint ODE coefficients are evaluated.
For a system with $n = 100$ state variables solved over $10{,}000$
time steps in double precision, the stored forward solution requires
$100 \times 10{,}000 \times 8$ bytes $= 8$~MB, modest.
For a climate model with $n = 10^6$ state variables and $10^6$
time steps, the requirement is $8 \times 10^{12}$~bytes $= 8$~TB;
infeasible to store in memory.

For large systems, two strategies address this:
\textbf{Checkpointing}\index{checkpointing!adjoint} stores the forward solution only at a subset
of time points and recomputes intermediate values as needed during
the backward pass. \textbf{Adjoint-of-checkpoint} schemes
(Griewank \& Walther~\cite{griewank2008evaluating}) achieve optimal trade-offs
between memory and recomputation cost.
For the models in this course, dense storage is always
adequate.

\subsection{Adaptive Step-Size Solvers}

As noted in Chapter~3, adaptive ODE solvers (such as \texttt{ode45})
introduce non-smoothness in the model output as a function of parameters:
the step-size controller makes discrete decisions that change discontinuously
with the parameter value. The continuous adjoint ODE from Chapter~6
assumes a smooth forward solution; when adaptive solvers introduce
discontinuities, the continuous adjoint may give inaccurate results.

The safest approach: use fixed-step-size integration (e.g., \texttt{ode4}
in MATLAB, a 4th-order Runge-Kutta) for the forward solve when
computing adjoint sensitivities. This eliminates the non-smoothness
and makes the continuous and discrete adjoints equivalent.
Alternatively, use tight tolerances (\texttt{RelTol} $\leq 10^{-8}$)
with adaptive solvers to minimize the non-smoothness.

\begin{selfcheck}
You compute adjoint sensitivities for the SIR model using
\texttt{ode45} with default tolerances for the forward solve,
then verify against centered finite differences.
The adjoint gives $\SI{\beta}{J} = 1.04$ while centered FD gives $0.87$.
Name two possible causes of this discrepancy, and describe
the diagnostic steps you would take to identify which cause applies.
\end{selfcheck}
\begin{scAnswer}
Possible causes: (1)~The adaptive step-size control in \texttt{ode45}
introduces non-smoothness that corrupts the continuous adjoint.
(2)~The finite-difference step size $h$ was chosen poorly (e.g.,
too small, causing catastrophic cancellation, or too large, causing
truncation error). Diagnostics: (a)~Recompute the adjoint with
fixed-step-size \texttt{ode4}: if the discrepancy disappears,
the solver is the culprit. (b)~Plot the U-curve for the FD estimate
to verify the step size is optimal; if the FD value changes
significantly with $h$, the FD is unreliable.
Always use finite differences as a sanity check for any adjoint
implementation; a discrepancy above $1\%$ warrants investigation.
\end{scAnswer}

\subsection{Discontinuities Break AD}\index{algorithmic differentiation!discontinuities}

AD requires the model to be differentiable at the nominal parameter values.
Common sources of non-differentiability in practice include \texttt{if} statements
whose conditions depend on parameters,\footnote{%
  \textit{If-statements are where smooth calculus goes to lose its temper.}
  When an \texttt{if} condition switches branches as a parameter is perturbed,
  the forward and reverse passes follow different computational paths,
  and the resulting ``derivative'' is the derivative of whichever branch
  happened to be active, which may be zero or undefined at the transition.}
\texttt{min} and \texttt{max} functions,
piecewise-defined model components such as quarantine thresholds or saturation
functions with hard cutoffs, and look-up tables interpolated without smoothing.
When these are present, AD will silently compute the derivative of one
branch, which may be zero or undefined at the branch point.
The diagnostic: compare AD results against centered finite differences.
If they disagree significantly for small $h$, suspect a discontinuity.
Replacing piecewise functions with smooth approximations (sigmoid
functions, smooth max/min) restores differentiability.

\begin{selfcheck}
A disease model contains the piecewise function:
$\beta(I) = \beta_0$ if $I < 0.1N$, else $\beta_0/2$.
This represents reduced contact when hospitalizations are high.
(a) Is this function differentiable with respect to $I$ everywhere?
(b) If you run reverse-mode AD through this function at $I = 0.09N$,
what derivative will AD report? Is this reliable?
(c) Propose a smooth approximation that would make the model
fully differentiable.
\end{selfcheck}
\begin{scAnswer}
(a) No. The function has a jump discontinuity at $I = 0.1N$,
so $\partial\beta/\partial I$ is undefined there.
(b) At $I = 0.09N$, the \texttt{if} condition is false ($I < 0.1N$),
so AD will compute the derivative of $\beta_0$ with respect to $I$,
which is zero. This is locally correct for $I$ strictly below the
threshold, but gives no information about the discontinuity itself,
and will be completely wrong for sensitivity with respect to $I$ near
the threshold. It is unreliable as a global SA result.
(c) Replace the step with a sigmoid:
$\beta(I) = \beta_0\bigl(1 - 0.5\,\sigma(k(I/N - 0.1))\bigr)$,
where $\sigma(x) = 1/(1+e^{-x})$ and $k > 0$ controls sharpness.
As $k \to \infty$ this recovers the discontinuous original;
a moderate $k$ (e.g., $k = 100$) approximates it closely while
remaining differentiable everywhere.
\end{scAnswer}

\unifyingidea{1} Every method in this chapter (forward mode,
reverse mode, symbolic differentiation, finite differences, the adjoint
ODE) computes the same object: Unifying Idea~1, the sensitivity index
defined in Chapter~2. The methods differ in computational cost,
accuracy, and applicability to different model types, but the result they
produce is always the same number: $\SI{p_j}{J} = (p_j/J)(dJ/dp_j)$.

%%------------------------------------------------------------
\section{MATLAB Laboratory: Four Methods, One Model}
\label{sec:ch7matlab}
%%------------------------------------------------------------

This laboratory computes the same sensitivity ($\partial J/\partial\beta$
for the SIR model with $J = \int_0^{90} I\,dt$, using all four methods
and compares their accuracy, speed, and ease of implementation.

\textbf{Part 1: Finite differences (Chapter~3 method).}
Use the \texttt{sensitivity\_jacobian} template from Chapter~3 with the
SIR model. Record the estimated $\SI{\beta}{J}$ and the computation time.

\textbf{Part 2: Analytic FSE (Chapter~4 method).}
Solve the 15-equation augmented system from Chapter~4 for the $\beta$
parameter alone. Extract $\partial I(t)/\partial\beta$ at all time points
and integrate numerically. Record the result and time.

\textbf{Part 3: Adjoint ODE (Chapter~6 method).}
Solve the adjoint system from Chapter~6. Use
formula~\eqref{eq:sensfml_vec} to compute $dJ/d\beta$.
Record the result and time.

\textbf{Part 4: MATLAB Symbolic Toolbox.}
Define the SIR ODE symbolically and compute the sensitivity
of the equilibrium $R_0 = k\beta/\gamma$ with respect to $\beta$
as a sanity check. Note: the Symbolic Toolbox computes sensitivities
of closed-form expressions, not of ODE solutions over time.

\begin{verbatim}
syms k beta gamma positive
R0   = k * beta / gamma;
SI_b = simplify( (beta/R0) * diff(R0, beta) )

%% For the time-dependent sensitivity via dlarray:
%% (requires Deep Learning Toolbox)
p_nom = [5, 0.06, 7, 10000];    %% [k, beta, tau, L]
p_dl  = dlarray(p_nom);
J_fn  = @(p) SIR_integral(p);   %% function returning integral of I
dJdp  = dlgradient(J_fn(p_dl), p_dl);
fprintf('AD sensitivity for beta: %.6f\n', extractdata(dJdp(2)));
\end{verbatim}

\textbf{Part 5: Comparison table.}
Fill in the following table from your measurements:

\begin{center}
\begin{tabular}{lcccc}
\toprule
Method & $\SI{\beta}{J}$ & Rel.\ error & Time (s) & Lines of new code \\
\midrule
Finite differences & & & & \\
Analytic FSE & & & & \\
Adjoint ODE & & & & \\
AD (\texttt{dlarray}) & & & & \\
\bottomrule
\end{tabular}
\end{center}

Discuss: which method is easiest to implement? Most accurate?
Fastest? Which would you use in practice for this model?

\textbf{Extension: introducing a discontinuity.}
Modify the SIR model to include a treatment threshold: if $I(t) > 0.1N$,
reduce $\tau$ by $20\%$ (faster treatment when infection is high).
Re-run all four methods. Which methods give consistent results?
Which methods silently produce wrong answers? How does the U-curve
(Chapter~3) change for finite differences near $\beta$?

%%------------------------------------------------------------
\paragraph{Connection to optimization.} The adjoint computed in this
chapter is the same object as the Lagrange multiplier in constrained
optimization and the costate variable in optimal control.
Appendix~\ref{app:optimization} develops this connection through
LP duality and Pontryagin's minimum principle.

%%------------------------------------------------------------
\section{Exercises}
%%------------------------------------------------------------

\begin{exercise}[\bloom{L3}]
\label{ex:ch7:annuity}
For the annuity function~\eqref{eq:annuity} at nominal values
$P = 1000$, $r = 0.05$, $n = 20$:
\begin{enumerate}[label=(\alph*)]
  \item Compute $A$ numerically.
  \item Compute all three sensitivities $\SI{P}{A}$, $\SI{r}{A}$, $\SI{n}{A}$
        analytically by differentiating~\eqref{eq:annuity}.
  \item Implement forward mode by hand for $q = r$ (as in Section~7.1)
        and verify your answer from (b).
  \item Implement reverse mode for all three parameters in one pass
        and verify all three answers.
  \item Count the number of arithmetic operations in each approach.
\end{enumerate}
\end{exercise}

\begin{exercise}[\bloom{L4}]
\label{ex:ch7:decisiontable}
For each of the following SA problems, select the most appropriate
method from Table~\ref{tab:masterSA} and justify in two sentences:
\begin{enumerate}[label=(\alph*)]
  \item A climate model with 800 parameters written in Fortran.
        Each model run takes 10~hours. You want sensitivities for
        2 response functionals.
  \item A pharmacokinetic model with 5 parameters and analytic
        closed-form solutions. You want sensitivities for 8 drug
        concentration time points.
  \item A traffic simulation with 40~parameters implemented as a
        compiled C++ executable. Parameter values can be set via
        a configuration file. Each run takes 3~minutes.
  \item A 10-parameter epidemic model written in MATLAB with
        \texttt{ode45}. You want one response functional.
        The model contains an \texttt{if} statement that triggers
        a quarantine intervention when $I > 0.05N$.
\end{enumerate}
\end{exercise}

\begin{exercise}[\bloom{L4}]
\label{ex:ch7:dlarray}
Using MATLAB's Deep Learning Toolbox \texttt{dlarray}:
\begin{enumerate}[label=(\alph*)]
  \item Wrap the function $R_0(k, \beta, \tau) = k\beta\tau$ as a
        differentiable computation and use \texttt{dlgradient} to compute
        all three partial derivatives at $k=5$, $\beta=0.06$, $\tau=7$.
  \item Compare with the analytic derivatives.
  \item Wrap the SIR equilibrium calculation (solve $\vec{f}(\vec{u}) = 0$
        numerically) and compute $\partial\bar{I}/\partial\beta$ at the
        endemic equilibrium. Compare with the FSE result from Chapter~4.
  \item What happens when you try to apply \texttt{dlgradient} through
        a call to \texttt{ode45}? Explain why, and what the remedy is.
\end{enumerate}
\end{exercise}

\begin{exercise}[\bloom{L5}]
\label{ex:ch7:workflow}
Design the complete SA workflow for a dengue fever transmission
model~\cite{chitnis2008determining} with the following properties:
$n = 8$ state variables (human and mosquito compartments),
$m = 14$ parameters (mortality rates, biting rates, transmission
probabilities, mean durations), $r = 2$ response functionals
(basic reproduction number and endemic human infection prevalence),
each ODE solve taking approximately 2~minutes.

Specify:
\begin{enumerate}[label=(\alph*)]
  \item Which SA method you would use and why.
  \item The expected total computation time.
  \item How you would validate the implementation.
  \item How you would handle the fact that two parameters
        ($\sigma_v$ and $\sigma_h$, the mosquito and human biting rates)
        are structurally correlated.
  \item What output format you would use to present results to
        a non-mathematical public health audience.
\end{enumerate}
\end{exercise}

\begin{exercise}[\bloom{L6}]
\label{ex:ch7:disagree}
A research group uses reverse-mode AD to compute adjoint sensitivities
for a mosquito population model with adaptive-step-size ODE integration.
They report that some sensitivity indices agree with finite differences
but others disagree by $10$--$20\%$, with no obvious pattern.
The discrepancies are largest for parameters that appear in the
mortality rate expressions.
\begin{enumerate}[label=(\alph*)]
  \item Identify the most likely technical cause of the discrepancy.
  \item Describe two diagnostic tests that would confirm or rule out
        this cause.
  \item Propose two remedies and explain the trade-offs between them.
  \item At what point would you conclude that the disagreement indicates
        a genuine SA result (the linearization is poor for these parameters)
        rather than a numerical artifact?
\end{enumerate}
\end{exercise}

%%------------------------------------------------------------

